% Generado por sistema/backend ReporteTex.java -- no editar a mano.
% Datos: Base de datos: 148 importado (titulados) / Base de datos: 53 importado (empleadores)
\subsubsection*{F4. Medio por el que obtuvo el trabajo}

Verificar si los titulados consiguen trabajo por todos los medios por igual o si hay medios que predominan.

\textit{Paso 1. Contar con la muestra y determinar su tamaño.} Población: trabajan en una organización; tamaño $n = 107$; opciones de la pregunta $k = 5$; $\alpha = 0.05$.

\begin{itemize}[nosep]
  \item $H_0$: Las respuestas de «Medio» se reparten por igual entre las 5 opciones (uniformidad).
  \item $H_1$: Las respuestas de «Medio» no se reparten por igual: hay opciones preferidas.
\end{itemize}

\textit{Paso 2. Calcular la frecuencia esperada.} Si las respuestas se repartieran por igual, cada opción tendría:
\[
FE = \frac{n}{k} = \frac{107}{5} = 21.40
\]

\textit{Paso 3. Calcular el estadístico chi-cuadrado.} Se cuenta la frecuencia observada $FO_i$ de cada opción y se calcula $\frac{(FO_i - FE)^2}{FE}$.

\begin{mitabla}{F4: prueba de las frecuencias, Medio}
\begin{tabular}{L{6cm}ccc}
\toprule
Opción & $FO_i$ & $FE$ & $\frac{(FO_i-FE)^2}{FE}$ \\
\midrule
Medios escritos/audiovisuales & 3 & 21.40 & $\frac{(3-21.40)^2}{21.40} = 15.8206$ \\[3pt]
Internet y redes sociales & 19 & 21.40 & $\frac{(19-21.40)^2}{21.40} = 0.2692$ \\[3pt]
Postulación competitiva & 44 & 21.40 & $\frac{(44-21.40)^2}{21.40} = 23.8673$ \\[3pt]
Recomendación de terceros & 32 & 21.40 & $\frac{(32-21.40)^2}{21.40} = 5.2505$ \\[3pt]
Otros & 9 & 21.40 & $\frac{(9-21.40)^2}{21.40} = 7.1850$ \\[3pt]
\midrule
Total & 107 & 107 & $x_0^2 = 52.3925$ \\
\bottomrule
\end{tabular}
\end{mitabla}

\[
x_0^2 = \sum_{i=1}^{5} \frac{(FO_i - FE)^2}{FE} = 52.3925
\]

\textit{Paso 4. Determinar el estadístico teórico.} $x^2_{\alpha,\,k-1} = x^2_{0.05,\,5-1} = x^2_{0.05,\,4} = 9.488$ (tabla chi-cuadrada).

\textit{Paso 5. Establecer la relación entre estadísticos.} $x_0^2 < x^2_{\alpha,\,k-1}$ $\rightarrow$ $52.3925 \ge 9.488$ $\rightarrow$ no.

\textit{Paso 6. Conclusión.}

\begin{center}
\begin{tabular}{l}
Si ($x_0^2 < x^2_{\alpha,\,k-1}$)\\
\quad ``Uniformidad''\\
Sino\\
\quad \fbox{No ``Uniformidad''}\\
FinSi
\end{tabular}
\end{center}

Se rechaza $H_0$: las respuestas no se reparten por igual, hay opciones preferidas. La más elegida es «Postulación competitiva» (44 de 107, 41.1\%).

\textit{Recomendación.} Hay medios que predominan: la Carrera debe apoyar la inserción laboral por esos medios (por ejemplo, bolsa de trabajo o red de contactos).

